\section{Comparación de aplicaciones web}

Según la información obtenida en el análisis anterior, podemos identificar una serie de funcionalidades comunes. No obstante, aunque las más populares son de nuestro interes \emph{(i.e. ya que nos indican lo que el usuario final puede esperar)}, a continuación nos centraremos en analizar las diferencias entre las aplicaciones:

Por un lado, en lo que respecta a las \textbf{utilidades presentes en las publicaciones de canciones} (véase Tabla~\ref{tab:comparativa_utilidades_canciones}), encontramos que la gran mayoría son ofrecidas por todas las aplicaciones. Si nos centramos en las utilidades que no son comunes en todas las páginas, podemos destacar las siguientes:

\begin{itemize}
    \item \textbf{Cambio de notación musical}, no disponible dada la popularidad de la notación inglesa en instrumentos de cuerda y al origen anglosajón de la mayoría de las aplicaciones. Solo LaCuerda ofrece la posibilidad de cambiar a notación latina, posiblemente por su origen hispanohablante.
    \item \textbf{Cambio de instrumento}, estando disponible en la mayoría de aplicaciones, salvo Ukulele Tabs que se centra exclusivamente en el ukelele.
    \item \textbf{Guía de rasgueos}, presente unicamente en Ukulele Tabs y Ultimate Guitar. Aunque no entraña dificultad técnica implementarla, la gran variedad de rasgueos posibles dificultar su inclusión en todas las aplicaciones.
\end{itemize}

\begin{table}[H]
    \centering
    \caption{Comparativa de utilidades en publicaciones por aplicación web.}\label{tab:comparativa_utilidades_canciones}
    \scriptsize
    \begin{tabular}{p{0.3\textwidth} c c c c}
        \hline
        \textbf{Utilidad} & \textbf{Ukulele Tabs} & \textbf{UkuTabs} & \textbf{Ultimate Guitar} & \textbf{LaCuerda} \\
        \hline
        Transportador & \ding{51} & \ding{51} & \ding{51} & \ding{51} \\
        Lista de acordes & \ding{51} & \ding{51} & \ding{51} & \ding{51} \\
        Autodesplazamiento & \ding{51} & \ding{51} & \ding{51} & \ding{51} \\
        Sistema de “me gusta” & \ding{51} & \ding{51} & \ding{51} & \ding{51} \\
        Sección de comentarios & \ding{51} & \ding{51} & \ding{51} & \ding{51} \\
        Impresión/Descarga & \ding{51} & \ding{51} & \ding{51} & \ding{51} \\
        Cambio de notación &  &  &  & \ding{51} \\
        Cambio de instrumento &  & \ding{51} & \ding{51} & \ding{51} \\
        Guardar canción & \ding{51} & \ding{51} & \ding{51} & \ding{51} \\
        Guía de rasgueos & \ding{51} &  & \ding{51} & \\
        \hline
    \end{tabular}
\end{table}

Por otro lado, en cuanto a las \textbf{funcionalidades generales de las aplicaciones web} (véase Tabla~\ref{tab:comparativa_funcionalidades}), podemos observar que la mayoría de aplicaciones ofrecen un conjunto similar de funcionalidades, siendo las únicas no comunes:

\begin{itemize}
    \item \textbf{Solicitud de nuevas canciones}, disponible en Ukulele Tabs, UkuTabs y LaCuerda. Ultimate Guitar no ofrece esta funcionalidad directamente, sino que lo hace a través de una sección de su foro.
    \item \textbf{Calculadora de escalas y notas}, presente en Ukulele Tabs y UkuTabs, posiblemente su ausencia se deba a la poca demanda de esta funcionalidad entre los usuarios de guitarra.
    \item \textbf{Diccionario de acordes}, disponible en Ukulele Tabs, UkuTabs y LaCuerda. Ultimate Guitar no ofrece esta funcionalidad debido a que ya sus publicaciones integran lista de acordes.
\end{itemize}

\begin{table}[H]
    \centering
    \caption{Comparativa de funcionalidades por aplicación web.}\label{tab:comparativa_funcionalidades}
    \scriptsize
    \begin{tabular}{p{0.3\textwidth} c c c c}
        \hline
        \textbf{Funcionalidad} & \textbf{Ukulele Tabs} & \textbf{UkuTabs} & \textbf{Ultimate Guitar} & \textbf{LaCuerda} \\
        \hline
        Biblioteca de canciones & \ding{51} & \ding{51} & \ding{51} & \ding{51} \\
        Publicaciones con letra y acordes & \ding{51} & \ding{51} & \ding{51} & \ding{51} \\
        Utilidades en publicaciones & \ding{51} & \ding{51} & \ding{51} & \ding{51} \\
        Impresión de publicaciones & \ding{51} & \ding{51} & \ding{51} & \ding{51} \\
        Creación y edición de publicaciones & \ding{51} & \ding{51} & \ding{51} & \ding{51} \\
        Puntuar publicaciones & \ding{51} & \ding{51} & \ding{51} & \ding{51} \\
        Comentar publicaciones & \ding{51} & \ding{51} & \ding{51} & \ding{51} \\
        Solicitar nuevas canciones & \ding{51} & \ding{51} &  & \ding{51} \\
        Cursos y guías & \ding{51} & \ding{51} & \ding{51} & \ding{51} \\
        Afinador & \ding{51} & \ding{51} & \ding{51} & \ding{51} \\
        Calculadora de escalas y notas & \ding{51} & \ding{51} &  &  \\
        Diccionario de acordes & \ding{51} & \ding{51} &  & \ding{51} \\
        \hline
    \end{tabular}
\end{table}

Podemos concluir que las utilidades y funcionalidades comunes son esenciales para el desarrollo de una aplicación web que siga el modelo de las páginas analizadas. Sin embargo, aquellas no comunes podrían ser consideradas como características adicionales que aporten valor añadido al producto final, diferenciándolo de la competencia.